% !TeX program = xelatex
\documentclass[11pt,a4paper]{article}

\usepackage{fontspec}
\setmainfont{DejaVu Serif}
\setmonofont{DejaVu Sans Mono}[Scale=0.82]

\usepackage{polyglossia}
\setdefaultlanguage{russian}
\setotherlanguage{english}

\usepackage{amsmath,amssymb}
\usepackage{booktabs}
\usepackage{geometry}
\geometry{left=22mm,right=20mm,top=20mm,bottom=20mm}
\usepackage{enumitem}
\usepackage{xcolor}
\usepackage{listings}
\usepackage[hidelinks]{hyperref}

% Без русских шаблонов переносов (texlive-lang-cyrillic) строки выходят
% за поля; эти настройки разрешают растягивать пробелы. При установленном
% пакете переносов можно убрать.
\tolerance=2000
\emergencystretch=3em
\hbadness=10000
\sloppy

\lstset{
  basicstyle=\ttfamily\footnotesize,
  breaklines=true,
  frame=single,
  framesep=3pt,
  rulecolor=\color{black!30},
  columns=fullflexible,
  keepspaces=true,
  extendedchars=true,
  commentstyle=\color{black!55}\itshape,
  morecomment=[l]{\#},
  aboveskip=6pt,
  belowskip=6pt
}

\newcommand{\bp}{\mathbf{p}}
\newcommand{\bq}{\mathbf{q}}
\newcommand{\bv}{\mathbf{v}}
\newcommand{\bs}{\mathbf{s}}
\newcommand{\bC}{\mathbf{C}}
\newcommand{\bR}{\mathbf{R}}
\newcommand{\bM}{\mathbf{M}}
\newcommand{\bB}{\mathbf{B}}
\newcommand{\bI}{\mathbf{I}}
\newcommand{\bQ}{\mathbf{Q}}
\newcommand{\bh}{\mathbf{h}}
\newcommand{\ba}{\mathbf{a}}
\newcommand{\bb}{\mathbf{b}}
\newcommand{\bc}{\mathbf{c}}
\newcommand{\bg}{\mathbf{g}}
\newcommand{\bzero}{\mathbf{0}}
\newcommand{\bPhi}{\boldsymbol{\Phi}}
\newcommand{\blam}{\boldsymbol{\lambda}}
\newcommand{\btau}{\boldsymbol{\tau}}
\newcommand{\bgam}{\boldsymbol{\gamma}}
\newcommand{\Ma}{\bM_{a}}
\newcommand{\tr}{^{\!\top}}
\newcommand{\code}[1]{\texttt{#1}}

\setlength{\parskip}{2pt}

\begin{document}

\begin{center}
{\large\bfseries Получение $\bM(\bq)$ и $\bh(\bq,\dot{\bq})$ из модели Flexia для планирования\\
оптимального по быстродействию движения}
\end{center}

\bigskip

\section*{Обозначения}

$\bp\in\mathbb R^{3n_b}$ --- абсолютные координаты положений тел,
$\bv=\dot{\bp}$;
$\blam\in\mathbb R^{m}$ --- множители Лагранжа;
$\bq\in\mathbb R^{d}$ --- суставные координаты, $d=3n_b-\operatorname{rank}\bC$;
$\bPhi(\bp)\in\mathbb R^{m}$ --- связи;
$\bC=\partial\bPhi/\partial\bp\in\mathbb R^{m\times3n_b}$;
$\Ma\in\mathbb R^{3n_b\times3n_b}$ --- матрица масс в абсолютных координатах.

Полный набор координат обозначен $\bp$, минимальный --- $\bq$; это разные
векторы. Знак при $\blam$ принят как в коде (вклад реакций в правой части).

\section{Задача}

Планировщик требует модель в явной форме
\begin{equation}
\btau=\bM(\bq)\ddot{\bq}+\bh(\bq,\dot{\bq}),
\qquad \btau_{\min}\le\btau\le\btau_{\max},
\label{eq:target}
\end{equation}
плюс путь $\bq(s)$, $s\in[0,1]$, дважды дифференцируемый. Такая форма нужна
потому, что после подстановки $\dot{\bq}=\bq'\dot s$ ограничение на моменты
становится линейным по $(\ddot s,\dot s^{2})$ и задача выпукла.

Flexia даёт систему в абсолютных координатах
\begin{equation}
\Ma\dot{\bv}=\bQ(t,\bp,\bv)+\bC\tr\blam,
\qquad
\bPhi(\bp)=\bzero,
\label{eq:dae}
\end{equation}
где $\bM_a$ постоянна и вырождена, размерность $96$ против $d=4$, управление
не выделено. Переход от~\eqref{eq:dae} к~\eqref{eq:target} --- редукция к
минимальным координатам, изложена ниже.

\paragraph{Сборка модели.} Для вычисления $\bM$ и $\bh$ модель собирается без
\code{PositionMotor2D} и \code{PositionLinearActuator2D}: в текущей
реализации привод --- кинематическая связь (\code{actuators.jl:18,44}), а не
источник обобщённой силы. Четыре привода в \code{parmech.jl} дали бы $d=0$.
Приводные углы должны оставаться свободными координатами, моменты --- входить
в правую часть как обобщённые силы.

\section{Редукция}

\paragraph{Скорости.} Дифференцируя $\bPhi(\bp)=\bzero$, получаем
$\bC\bv=\bzero$. Разделим индексы на независимые $\mathcal I$ ($d$ штук,
угловые координаты кривошипов) и зависимые $\mathcal D$ ($m$ штук). Тогда
$\bC_{\mathcal I}\bv_{\mathcal I}+\bC_{\mathcal D}\bv_{\mathcal D}=\bzero$,
откуда при $\bq:=\bp_{\mathcal I}$
\begin{equation}
\bv=\bR\dot{\bq},
\qquad
\bR_{\mathcal I,:}=\bI_d,
\qquad
\bR_{\mathcal D,:}=-\bC_{\mathcal D}^{-1}\bC_{\mathcal I},
\qquad
\bC\bR=\bzero.
\label{eq:R}
\end{equation}
Столбцы $\bR$ --- базис $\ker\bC$, то есть касательное пространство к
многообразию $\bPhi=\bzero$.

\paragraph{Ускорения.} Второе дифференцирование даёт $\bC\dot{\bv}=\bgam$,
$\bgam=-\dot{\bC}\bv$, откуда
\begin{equation}
\dot{\bv}=\bR\ddot{\bq}+\bs,
\qquad
\bs_{\mathcal I}=\bzero,
\qquad
\bs_{\mathcal D}=\bC_{\mathcal D}^{-1}\bgam.
\label{eq:s}
\end{equation}
Сопоставление с дифференцированием~\eqref{eq:R} даёт $\bs=\dot{\bR}\dot{\bq}$;
явное дифференцирование $\bR$ не требуется. Величина $\bgam$ есть производная
по направлению и берётся автодифференцированием:
\begin{equation}
\bgam=-\left.\frac{d}{d\varepsilon}\right|_{\varepsilon=0}\bC(\bp+\varepsilon\bv)\,\bv .
\end{equation}

\paragraph{Проекция.} Подставляя~\eqref{eq:s} в~\eqref{eq:dae} с добавлением
моментов приводов и умножая слева на $\bR\tr$, получаем
$\bR\tr\bC\tr=(\bC\bR)\tr=\bzero$ (реакции идеальных связей ортогональны
допустимым движениям), и множители исчезают:
\begin{equation}
\bM(\bq)\ddot{\bq}+\bh(\bq,\dot{\bq})=\bR\tr\bB\btau,
\qquad
\bM=\bR\tr\Ma\bR,
\qquad
\bh=\bR\tr\Ma\bs-\bR\tr\bQ.
\label{eq:Mh}
\end{equation}
Матрица $\bM$ симметрична и положительно определена при
$\operatorname{rank}\bR=d$; это же следует из
$\tfrac12\bv\tr\Ma\bv=\tfrac12\dot{\bq}\tr\bM\dot{\bq}$.

Столбец $\bB$ для момента $\tau_k$ содержит $\pm1$ в угловых строках двух тел
шарнира. Если привод соединяет закреплённое основание с телом, угол которого
выбран независимой координатой, то $\bR\tr\bB=\bI_d$ и
$\btau=\bM\ddot{\bq}+\bh$; иначе произведение считается явно.

\section{Извлечение исходных данных из Flexia}

Строки множителей в правой части содержат в точности $\bPhi$ (присваивание, а
не накопление: \code{joints.jl:68–70, 157–161, 391–395}), поэтому $\bC$
доступна как блок якобиана (\code{system.jl:117}). Вычитание уставок
(\code{system.jl:112}) на производную не влияет.

\begin{lstlisting}
lm_rows  = collect((last_body_dof(sys)+1):last_lm_dof(sys))
pos_cols = vcat([get_body_position_dofs(sys, b) for b in sys.bodies]...)
vel_cols = vcat([get_body_velocity_dofs(sys, b) for b in sys.bodies]...)

C  = sys.jacobian(state)[lm_rows, pos_cols]     # m x 3n_b
Ma = sys.mass[vel_cols, vel_cols]               # 3n_b x 3n_b, постоянная

st = copy(state); st[lm_rows] .= 0.0            # вклады связей линейны по L
Q  = sys.rhs(st)[vel_cols]                      # чистые внешние силы
\end{lstlisting}

Порядок индексов в \code{pos\_cols} задаёт нумерацию координат для всех
матриц и должен совпадать с порядком \code{vel\_cols} и с разбиением
$\mathcal I/\mathcal D$. Несогласованность даёт результат, сохраняющий
симметричность и положительную определённость $\bM$, то есть не выявляемый
стандартными проверками.

Прямой путь через \code{sys.kinematic\_constrains!} невозможен: сигнатуры
\code{Vector\{Float64\}} в \code{compute\_kinematic\_residual!}
(\code{joints.jl:72, 172, 348}) не пропускают дуальные числа. Замена на
\code{AbstractVector\{<:Real\}} сделала бы обходной путь ненужным.

Целесообразно оформить всё перечисленное отдельным модулем
\code{src/reduction.jl} со структурой, фиксирующей списки индексов, и
функциями \code{constraint\_jacobian}, \code{nullspace\_basis},
\code{mass\_matrix}, \code{bias\_vector}, \code{assemble\_configuration!}.
Дополнительно нужен тип \code{TorqueActuator2D <: AbstractForce2D} --- привод
как источник момента (по устройству аналогичен \code{TorsionalSpring},
\code{forces.jl:25–53}, \code{number\_of\_dofs} равен нулю). Сейчас моменты
приходится подавать через поле \code{forces} тел.

\section{Путь и форма для планировщика}

Путь $\bq(s)$ задаётся либо напрямую сплайном по узлам в суставных
координатах, либо через обратную задачу: для каждого $s_k$ решается
$[\bPhi(\bp);\ \bg(\bp)-\bg_{\text{зад}}(s_k)]=\bzero$ методом Ньютона со
стартом от предыдущего узла, из решения берутся компоненты $\mathcal I$.
Второй способ требует функции сборки конфигурации; \code{static\_solver!}
для этого не годится.

Подстановка $\dot{\bq}=\bq'\dot s$, $\ddot{\bq}=\bq'\ddot s+\bq''\dot s^{2}$
в~\eqref{eq:Mh} даёт
\begin{equation}
\btau(s)=\ba(s)\ddot s+\bb(s)\dot s^{2}+\bc(s),
\end{equation}
\begin{equation}
\ba=\bM(\bq)\bq',
\qquad
\bc=\bh(\bq,\bzero),
\qquad
\bb=\bM(\bq)\bq''+\bh(\bq,\bq')-\bc .
\end{equation}
Использовано то, что $\bh$ есть сумма однородного второй степени по скорости
слагаемого и позиционного. Матрицу Кориолиса отдельно вычислять не нужно:
все коэффициенты получаются тремя вызовами $\bM$ и $\bh$.

Разложение верно, только если $\bQ$ не зависит от скоростей. Демпфирование
(поле \code{damping} в \code{TorsionalSpring}) добавляет член, линейный по
$\dot s$, что нарушает выпуклость. В \code{parmech.jl} диссипации нет.

\end{document}
